\documentclass[12pt]{article}


\usepackage{bm}		
\usepackage{float}
\usepackage{array}
\usepackage{xcolor}
\usepackage{amsthm}
\usepackage{amsmath}
\usepackage{graphicx}
\usepackage{listings}
\usepackage{setspace}
\usepackage{enumitem}
\usepackage{amsfonts}	
\usepackage{mathtools}
\usepackage{subcaption}
\usepackage[utf8]{inputenc}
\usepackage[font=scriptsize,labelfont=bf]{caption}
\usepackage[inner=2.5cm,outer=2.5cm,top=2.5cm,bottom=2.5cm]{geometry}	

\lstset{
   language=C,
   backgroundcolor=\color{white},
   extendedchars=true,
   basicstyle=\footnotesize\ttfamily,
   showstringspaces=false,
   showspaces=false,
   numbers=left,
   numberstyle=\footnotesize,
   numbersep=9pt,
   tabsize=4,
   breaklines=true,
   showtabs=false,
   captionpos=b,
   frame=single,
    %===========================================================
    framesep=3pt,%expand outward.
    framerule=0.4pt,%expand outward.
    xleftmargin=3.4pt,%make the frame fits in the text area. 
    xrightmargin=3.4pt%make the frame fits in the text area.
    %=========================================================== 
}

\title{Projet compilation}
\author{Gabriel VALDÉS-ALONZO}
\date{3 mai 2025}

\begin{document}
\maketitle
\clearpage

\section{Introduction}
\section{Structure du code}
\subsection{Structure du projet}
\subsection{Structure du dossier}
Le projet est structuré de manière d'avoir une repartition logique des différents composant. La première division est au niveau des dossiers dans le projet, qui sont organisés de la manière suivante :
\begin{verbatim}
  Code
  |_ build/
  |_ doc/
  |_ files/
  |_ out/
  |_ src/
  |_ Makefile
\end{verbatim}

\noindent où les différents dossiers ont chacun sa fonction:

\begin{itemize}
  \item \texttt{build} : dossier reservé pour tous les sorties du fichier \texttt{Makefile}. Cela concerne tous les fichiers intermédiaires \texttt{.o} et aussi l'executable après la compilation.
  \item \texttt{doc} : dossier contenant ce documentation et les fichier tex utilisés pour le générer.
  \item \texttt{files} : dossier reservé pour les fichiers d'entré du code, particulièrement le fichier \texttt{.txt} avec l'expression régulier à analyser. 
  \item \texttt{out} : dossier reservé pour tous les sorties du code. Dans ce projet en particulier les sorties sont les fichiers \texttt{.dot} qui représentent les différents objets du code (les automates et l'arbre syntaxique), les images générés par \textit{graphviz} et le fichier \texttt{.c} generé à partir de l'automate minimale.
  \item \texttt{src} : dossier reservé pour le code source, c'est-à-dire tous les fichiers \texttt{.c} et \texttt{.h} avec les fonctions et codes implémentent les automates et le code de sortie.
\end{itemize}

Finalement, a la racine du projet on a simplement un fichier \texttt{Makefile}, qui est là pour simplifier la compilation du code. Le fichier peut être appellé avec la comande shell \texttt{make}. Si l'on veut refaire la compilation on peut appeler avant la commande \texttt{make clean}, commande qui appele simplement \texttt{rm -rf build}.

\subsection{Structure des fichiers}
\section{Documentation des fonctions}

\subsection{Main functions}
\subsection{aux\_structs}
\subsection{f\_aux\_automata}
\subsection{f\_aux\_tree}

\end{document}